% Rozdziały zaczynamy od nowej strony.
\subsection{Graphs Neural Networks}

% Graph Neural Networks (GNNs) are types of models that use on how the specific node
% is connected to other nodes in the graph. During training every node is represented 
% by its features and neighbors features that allows not only represent the node based 
% on its own characteristics like it is in traditional neural networks, but also represent
% the node based on its context. This mechanism is called message passing and allows to 
% aggregate information from direct neigbours as well further ones depending on assumed
% number of layers in network. There are multiple types of GNNs that differ from Each
% other based on how the message from neighbors is aggregated. In this Section we 
% review couple of them that were applied in this work.

Graph neural networks (GNNs) are a class of models that leverage how a given node is connected 
to other nodes in the graph. During training, every node is represented not only by its own 
features, as in traditional neural networks, but also by the features of its neighbours.
What it is new in GNNs comparing to previous graph embedding methods is that the representation
is not fixed and it is built based on aggregation, so it is not necessary to have all nodes in the graph 
during training. 
%fixed representation for every node observed during training, making it impossible to 
%generate embeddings for nodes that did not appear in the training graph. 
%This allows the node's representation to capture its surrounding context in addition to its own 
%characteristics. 
This mechanism, known as message passing, aggregates information from direct 
neighbours as well as from more distant ones, depending on the number of layers assumed in the 
network. GNN  builds a node's representation at each layer by combining its own representation from
the previous layer with the aggregated representations of its neighbors, also from the 
previous laye
Formally, the representation $h_v^{(K)}$ of node $v$ depends on
the features of all nodes within its $K$-hop neighborhood.
Since each layer $k$ updates
\begin{equation}
h_v^{(k)} = \text{UPDATE}\Bigl(h_v^{(k-1)},\ \text{AGGREGATE}\bigl(\{h_u^{(k-1)} : u \in \mathcal{N}(v)\}\bigr)\Bigr),
\end{equation}
so that information from a node $u$ at graph distance $d$ from $v$ only reaches $v$ once $K \geq d$.
Different types of GNNs differ primarily in how they aggregate messages from their 
neighbors. In this section, we review several such architectures that were applied in this work.

--dodanie informacji o wagach modelu i jak tworzona reprezentacja / embedding 

\subsubsection{GraphSAGE}

  GraphSAGE \cite{hamilton2017inductive} was introduced in response to the need for inductive 
  and scalable representation learning on graphs. As introduced in this section introduction, 
  GraphSAGE instantiates the general message-passing scheme using one of three 
  aggregation functions — mean, LSTM, or pooling — so that the same learned functions can later 
  be applied to any node.
  What is specific to GraphSAGE is that the node's own representation and the aggregated 
  neighbor representation are concatenated into single and longer vector, 
  and the resulting vector is L2-normalized after each layer.
  Additionally, GraphSAGE addresses the scalability of training by sampling a fixed number of 
  neighbors for each node, not full neighbourhood,
  and updating the node's representation based on this sampled subset.

\subsubsection{Relational Graph Convolutional Network}

Relational Graph Convolutional Networks (R-GCNs) \cite{schlichtkrull2018modeling}
extend Graph Convolutional Networks (GCNs) to multi-relational graphs, in
which edges represent different types of relationships. Unlike a
standard GCN, which applies the same transformation to messages from all
neighbours, R-GCN applies a relation-specific transformation. 

For a node $v$, the representation at layer $l+1$ is computed as

\begin{equation}
\label{eq:rgcn}
h_v^{(l+1)} =
\sigma \left(
\sum_{r \in \mathcal{R}}
\sum_{u \in \mathcal{N}_r(v)}
\frac{1}{c_{v,r}}
W_r^{(l)} h_u^{(l)}
+
W_0^{(l)} h_v^{(l)}
\right),
\end{equation}

where $\mathcal{R}$ is the set of relation types,
$\mathcal{N}_r(v)$ denotes the set of neighbours connected to node $v$
through relation $r$, $W_r^{(l)}$ is a trainable weight matrix specific to
relation $r$, and $c_{v,r}$ is a normalisation constant. The term
$W_0^{(l)} h_v^{(l)}$ represents a self-loop transformation, which preserves
the node's own representation in addition to messages received from its
neighbours. Finally, $\sigma$ denotes a non-linear activation function.

When the number of relation types is large, assigning a separate full
weight matrix to each relation can lead to a large number of parameters and overfitting.
To address this issue, R-GCN introduces basis decomposition, in which each
relation-specific matrix is expressed as a linear combination of $B$ shared
basis matrices:

\begin{equation}
\label{eq:rgcn-basis}
W_r^{(l)} =
\sum_{b=1}^{B} a_{rb}^{(l)} V_b^{(l)},
\end{equation}

where $V_b^{(l)} \in
\mathbb{R}^{d_{\mathrm{out}} \times d_{\mathrm{in}}}$ is the $b$-th shared
basis matrix and $a_{rb}^{(l)}$ is a trainable coefficient specifying the
contribution of this basis to relation $r$. 

\subsubsection{Transformed-based GNN}

Transformer-based graph neural networks adapt the self-attention mechanism
introduced in the Transformer architecture (\textcite{vaswani2017attention}) to
graph-structured data.
Instead of using a fixed type of aggregation of neighbours information,
the architecture uses an attention-based message passing mechanism that allows 
to recognize the importantance of neighbouring nodes.
For each node $i$, its input representation $\mathbf{h}_i$ is designed by a projections:
query, key and value:
\begin{equation}
\mathbf{q}_i = \mathbf{W}_{Q}\mathbf{h}_i, \qquad
\mathbf{k}_i = \mathbf{W}_{K}\mathbf{h}_i, \qquad
\mathbf{v}_i = \mathbf{W}_{V}\mathbf{h}_i,
\end{equation}
where $\mathbf{W}_{Q}$, $\mathbf{W}_{K}$, and $\mathbf{W}_{V}$ are learnable
weight matrices.
Then for an edge from a source node $j$ to a target node $i$, the weight of attention 
(attention coefficient) is calculated based on equation:
\begin{equation}
    e_{ij} =
    \frac{\mathbf{q}_{i}^{\top}\mathbf{k}_{j}}{\sqrt{d_k}},
    \end{equation}
    where $d_k$ denotes the dimensionality of the key vectors.

The attention coefficients are normalised across all neighbours $\mathcal{N}(i)$ of the target
    node using the softmax function:
    \begin{equation}
    \alpha_{ij} =
    \frac{\exp(e_{ij})}
    {\sum_{r \in \mathcal{N}(i)} \exp(e_{ir})}.
    \end{equation}
The node's representation is updated by aggregating value projections of
its neighbours, weighted by the learned attention coefficients:
\begin{equation}
\mathbf{h}'_i =
\sum_{j \in \mathcal{N}(i)}
\alpha_{ij}\mathbf{v}_{j}.
\end{equation}

To capture multiple complementary patterns of neighbourhood interactions, Graph
Transformers commonly use multi-head attention. For each attention head
$m \in \{1, \ldots, H\}$, independent projection matrices
$\mathbf{W}_{Q}^{(m)}$, $\mathbf{W}_{K}^{(m)}$, and
$\mathbf{W}_{V}^{(m)}$ are learned. The output of head $m$ for node $i$ is:
\begin{equation}
\mathbf{h}_i^{\prime(m)} =
\sum_{j \in \mathcal{N}(i)}
\alpha_{ij}^{(m)}\mathbf{v}_{j}^{(m)}.
\end{equation}
The outputs of all $H$ heads are combined by concatenating or averaging them.
This mechanism provides a trainable aggregation methods, not defined by default like in previously mentioned
architectures.

All architectures mentioned in this section can be applied to heterogeneous
graphs through an appropriate configuration of the \texttt{HeteroConv} layer in
PyTorch Geometric (PyG). \texttt{HeteroConv} is a generic wrapper that applies
a separately specified message-passing operator to each edge type, i.e., each
meta-relation between a source node type and a target node type
\textcite{wojcik2026grafowe}.

For attention-based message passing in heterogeneous graphs, a dedicated
architecture called the Heterogeneous Graph Transformer (HGT) was proposed 
 in \textcite{hu2020heterogeneous}.
HGT makes attention scores and message transformations dependent on the
types of source and target nodes and on the edge relation type.

For an edge from a source node $s$ to a target node $t$, HGT represents its
semantics as a meta-relation
$(\tau_s, \phi(e), \tau_t)$, where $\tau_s$ and $\tau_t$ denote the source
and target node types, respectively, and $\phi(e)$ denotes the edge type.
The target node produces a query, whereas the source node produces a key and
a value. Relation-specific transformations affect both the attention score and
the message transferred from $s$ to $t$.

Consequently, HGT can assign different importance to neighbours depending not
only on their features, but also on the semantic type of their connection to
the target node. The weighted incoming messages are then aggregated and
transformed using parameters specific to the target node type. 

%\subsubsection{Depth of GNNs}
% Residual connections switch the effective per-layer propagation operator from the
% (nilpotent, on a DAG without self-loops) 
% adjacency operator , which is the standard formalization of the over-smoothing regime and is the mechanism 
% through which multi-layer signal survives on a directed acyclic structure.

% PairNorm (Zhao et Akoglu, 2020) re-centers and re-scales node representations after each layer to keep the pairwise distance between representations from collapsing, without preventing the underlying mixing.

% DropEdge (Rong et al., 2020) randomly removes a fraction of edges independently in every layer during training, slowing the rate of information mixing and acting as a structural regularizer.

% Jumping Knowledge (Xu et al., 2018) aggregates the outputs of all layers (via concatenation, cat, or elementwise max) rather than only the last one, letting a node fall back on a shallower, less-smoothed representation when the deepest layer has become uninformative.

% Weisfeiler-Lehman --- two nodes with identical initial features and structurally equivalent neighborhoods 
% cannot be assigned distinct representations by any number of convolutional layers.

% leave-one-out